\subsection{Lattice Points in Circles}
\label{subsections:examples:example3}



\begin{exercise}{
    A point $(m, n) \in \R^2$ is a lattice point if $m, n \in \Z$. %
    Prove that the number of lattice points inside any circle centered at %
    the origin equals $4k + 1$ for some $k \in \Z$.
}
\end{exercise}



\begin{proof}{
    This statement can be proven by %
    \hyperref[subsections:techniques:universal:cases]{division into cases}.

    \begin{panel}[Setup.]{
        Let $r \in \R$ such that $r > 0$ is the radius of a circle %
        centered at the origin, $(0, 0)$. %
        Let $(m, n)$ be a lattice point inside the circle. %
        Let $d$ denote the distance from $(0, 0)$ to $(m, n)$. %
        By the Euclidean distance formula, %
        \begin{align*}{
            d &= \sqrt{(m - 0)^2 + (n - 0)^2}\\
              &= \sqrt{m^2 + n^2}\\
              &< r.
        }
        \end{align*}
        Thus, $m^2 + n^2 < r^2$. %
        Since $m$ and $n$ are integers, each is either $0$ or nonzero. %
        Therefore, we consider three cases: both are $0$, %
        exactly one is $0$, or neither is $0$.
    }
    \end{panel}

    \begin{panel}[Case 1: Suppose that $m = 0$ and $n = 0$.]{
        Observe that $0^2 + 0^2 = 0 < r^2$. %
        Therefore, the origin $(0, 0)$ is inside the circle, %
        giving exactly $1$ lattice point.
    }
    \end{panel}

    \begin{panel}[Case 2: Suppose that exactly one of $m$ and $n$ is $0$.]{
        Let $t \in \Z$ such that $t > 0$ and $t$ is the absolute value of %
        the nonzero coordinate of $(m, n)$. %
        Then $(m, n)$ is one of the points %
        \[
            (t, 0), \quad (-t, 0), \quad (0, t), \quad (0, -t).
        \]
        These are four distinct lattice points, since $t \neq 0$. %
        Since the squares of the coordinates of each point sum to %
        $t^2 = m^2 + n^2 < r^2$, all four points are inside the circle. %
        Since $t$ is determined by the point, distinct values of $t$ give %
        disjoint sets. %
        Therefore, the lattice points in this case are partitioned into %
        sets of $4$.
    }
    \end{panel}

    \begin{panel}[Case 3: Suppose that $m \neq 0$ and $n \neq 0$.]{
        Let $s = |m|$ and $t = |n|$. %
        Then $s, t \in \Z$ with $s > 0$ and $t > 0$, %
        and $(m, n)$ is one of the points %
        \[
            (s, t), \quad (-s, t), \quad (s, -t), \quad (-s, -t).
        \]
        These are four distinct lattice points, since $s \neq 0$ and %
        $t \neq 0$. %
        Since $(\pm s)^2 + (\pm t)^2 = m^2 + n^2 < r^2$, %
        all four points are inside the circle. %
        Since the pair $(s, t)$ is determined by the point, distinct pairs %
        give disjoint sets. %
        Therefore, the lattice points in this case are partitioned into %
        sets of $4$.
    }
    \end{panel}

    \begin{panel}[Combining the cases.]{
        Since $m^2 + n^2 < r^2$ implies $|m| < r$ and $|n| < r$, %
        only finitely many lattice points are inside the circle. %
        Since the lattice points in Cases 2 and 3 are partitioned into sets %
        of $4$, there exists $k \in \Z$ with $k \geq 0$ such that they %
        total $4k$ points. %
        Adding the $1$ point from Case 1, %
        the circle contains $4k + 1$ lattice points.
    }
    \end{panel}

    Therefore, the statement holds.
}
\end{proof}